% from ega4-18-fr.tex
\oldpage[IV-4]{186}

irréductible de $Y$ contenant $y$; comme un tel point $z$ est contenu dans tout voisinage de $y$ dans $X$, $\mathscr J_z/\mathscr J_z^2$ est un $A$-module libre de rang $n=\operatorname{codim}_y^*(Y,X)$; tout revient à voir que $\dim(A)=n$. Or, puisque $z$ est point maximal du sous-préschéma $Y$ défini par $\mathscr J$, on a $\dim(\sh{O}_{X,z}/\mathscr J_z)=0$, donc $\mathscr J_z$ est un \emph{idéal de définition} de l'anneau local noethérien $A$; en outre, l'hypothèse que $\mathscr J$ est quasi-régulier entraîne que $\mathscr J_z^m/\mathscr J_z^{m+1}$ est un $(A/\mathscr J_z)$-module libre de rang $\binom{m+n-1}{n-1}$; si $r$ est la longueur de l'anneau artinien $A/\mathscr J_z$, la longueur de $\mathscr J_z^m/\mathscr J_z^{m+1}$ est donc $r\binom{m+n-1}{n-1}$; le polynôme de Hilbert-Samuel de $A$ pour la filtration $\mathscr J_z$-préadique est donc bien de degré $n$, ce qui prouve la proposition ($0$, 16.2.3).

En vertu de (19.1.4), on dira désormais « codimension » au lieu de « codimension transversale » lorsqu'il s'agira d'un sous-préschéma $Y$ régulièrement immergé dans $X$ (\emph{même lorsque $X$ n'est pas localement noethérien}).

\begin{proposition}[19.1.5]
\label{IV.19.1.5}
\begin{enumerate}
  \item[(i)] Pour qu'une immersion $j:Y\to X$ soit ouverte, il faut et il suffit qu'elle soit régulière et de codimension $0$ en tout point.
  \item[(ii)] Soient $f:Y\to X$ une immersion régulière (resp. quasi-régulière), $g:X'\to X$ un morphisme plat, $Y'=Y\times_X X'$; alors le morphisme $f'=f_{(X')}:Y'\to X'$ est une immersion régulière (resp. quasi-régulière), et pour tout $y'\in Y'$, si $y$ est la projection de $y'$ dans $Y$, la codimension de $f'$ au point $y'$ est égale à la codimension de $f$ au point $y$. Inversement, si $f$ est une immersion et $g:X'\to X$ un morphisme fidèlement plat et quasi-compact, alors, si $f'$ est une immersion quasi-régulière, il en est de même de $f$.
  \item[(iii)] Si $f:Y\to X$ et $g:Z\to Y$ sont des immersions régulières, alors $f\circ g:Z\to X$ est une immersion régulière, et pour tout $z\in Z$, la codimension de $f\circ g$ au point $z$ est la somme de la codimension de $g$ au point $z$ et de la codimension de $f$ au point $g(z)$. De plus, la suite d'homomorphismes canoniques (16.2.7.1)
\begin{equation}
  0\longrightarrow g^*(\sh{N}_{Y/X})\longrightarrow\sh{N}_{Z/X}\longrightarrow\sh{N}_{Z/Y}\longrightarrow0
  \tag{19.1.5.1}\label{IV.19.1.5.1}
\end{equation}
est exacte, et pour tout $z\in Z$, il existe un voisinage de $z$ dans $Z$ dans lequel les restrictions des homomorphismes de cette suite forment une suite scindée.
  \item[(iv)] Soient $X$ un préschéma localement noethérien, $f:Y\to X$ et $g:Z\to Y$ deux immersions, $z$ un point de $Z$. Les conditions suivantes sont équivalentes:
  \begin{enumerate}
    \item[a)] $g$ est une immersion régulière au point $z$ et $f$ une immersion régulière au point $g(z)$.
    \item[b)] $g$ et $f\circ g$ sont des immersions régulières au point $z$.
    \item[c)] $f\circ g$ est une immersion régulière au point $z$, $f$ est une immersion régulière au point $g(z)$, et la suite d'homomorphismes canoniques (19.1.5.1), restreinte à un voisinage ouvert assez petit de $z$ dans $Z$, est exacte et scindée.
  \end{enumerate}
  \item[(v)] Soient $X$ un préschéma localement noethérien, $f:Y\to X$, $g:Z\to Y$ deux immersions, $z$ un point de $Z$; supposons que $y=g(z)$ soit un point régulier de $Y$ et $x=f(g(z))$ un point régulier

